\newpage
\section{ANTECEDENTES}

El monitoreo de condición de herramientas de corte constituye una estrategia relevante para mejorar la calidad de manufactura. Permite disminuir la generación de piezas defectuosas y reducir paradas no planificadas en máquinas CNC. Dentro de las técnicas indirectas de monitoreo, el análisis de señales eléctricas de los accionamientos representa una alternativa de bajo impacto. Esto permite inferir cambios en la carga de corte sin necesidad de instalar sensores invasivos directamente sobre la herramienta de corte.

En este enfoque, las variaciones de desgaste, pérdida de filo o rotura modifican las fuerzas de mecanizado. En consecuencia, se altera la corriente consumida por el motor del husillo o por los motores de avance. Un antecedente temprano y fundamental de este enfoque es el trabajo de \cite{romero2003}, quienes propusieron detectar la rotura de herramientas en fresadoras CNC. Lograron esto mediante el análisis continuo de la corriente de los accionamientos durante el mecanizado.

Su propuesta permitió identificar eventos asociados a la rotura sin recurrir a dinamómetros ni a sensores de fuerza de alto costo. Además, evitó realizar modificaciones importantes en la estructura mecánica de la máquina, lo que consolidó el uso de señales internas. A este enfoque indirecto y pragmático se le denomina frecuentemente en la literatura como monitoreo \textit{sensorless} o de instrumentación mínima.

La línea de investigación fue ampliada por \cite{trejo2010}, quien estudió el monitoreo y modelado fenomenológico del desgaste de herramienta. Las pruebas se realizaron bajo condiciones variables de corte en torno CNC, integrando señales de corriente y vibración. Se utilizó un sensor inteligente basado en FPGA con el propósito fundamental de estimar cuantitativamente la evolución del desgaste de flanco a lo largo del proceso.

Este antecedente es importante porque no se limita a identificar una falla súbita aislada durante la operación. Al contrario, aborda la evolución progresiva del desgaste y su comportamiento en condiciones de corte dinámicas y variables. Posteriormente, \cite{sevilla2011} validaron un método de detección de rotura de herramienta en fresado de alta velocidad analizando las señales de corriente del motor de avance.

El estudio de \cite{sevilla2011} utilizó la transformada wavelet discreta para extraer información transitoria de la señal analizada. Además, aplicó herramientas estadísticas robustas para discriminar eficazmente condiciones normales y anormales de mecanizado. Este valioso aporte respalda que las señales de corriente pueden contener patrones sumamente útiles para la detección de eventos súbitos aplicando procesamiento tiempo-frecuencia.

Como complemento alternativo, \cite{ramirez2018} desarrollaron un sensor inteligente para detección de rotura en fresado utilizando termografía infrarroja. Estas pruebas experimentales se realizaron rigurosamente bajo condiciones de corte seco y también utilizando abundante refrigerante. Sin embargo, la termografía requiere forzosamente una línea de visibilidad directa hacia la zona de corte en el interior de la máquina.

Este requerimiento visual puede verse afectado negativamente por la presencia de viruta, refrigerante o diversas restricciones físicas de la máquina CNC. Por ello, el monitoreo basado en el análisis de corriente mantiene una ventaja sumamente relevante y pragmática frente a los sistemas ópticos. Destaca de manera especial cuando se busca una implementación robusta de fácil instalación y de muy bajo costo computacional.

La vigencia del monitoreo no invasivo por corriente se confirma fehacientemente con investigaciones recientes que incorporan aprendizaje profundo. \cite{lai2025} propusieron un sistema eficiente basado en señales de corriente del husillo principal. Dichas señales fueron obtenidas con un sensor no invasivo de tipo SCT013 instalado externamente alrededor del cableado, sin intervenir para nada el controlador.

Los autores analizaron las señales mediante la transformada rápida de Fourier y compararon diversos modelos de redes neuronales artificiales (ANN, DNN, CNN). Su objetivo primordial fue clasificar eventos de rotura, logrando detectar graves anomalías con una extraordinaria precisión y destacable rapidez. La utilización de una pinza SCT013 es particularmente pertinente para instrumentación de muy bajo impacto y fácil despliegue.

El sensor SCT013 permite implementar un sistema de adquisición de datos totalmente reversible y seguro para maquinaria heredada sin conectividad nativa. No obstante, el diseño meticuloso del sistema debe considerar rigurosamente el rango nominal de la corriente de trabajo y la frecuencia de muestreo adecuada. También requiere implementar de manera obligatoria un óptimo acondicionamiento de señal para lograr rechazar el elevado ruido eléctrico industrial en planta.

Por otra parte, \cite{wang2024} propusieron una innovadora metodología de predicción basada estrictamente en la fusión de información multisensorial. El estudio científico integra conjuntamente complejas señales de corriente, niveles de vibración mecánica y medición directa de la fuerza de corte. Esto demuestra fehacientemente que las variables eléctricas y mecánicas siempre logran aportar información complementaria para un diagnóstico holístico.

Una limitación frecuente de los modelos tradicionales es su drástica disminución de desempeño general cuando cambian drásticamente las condiciones operativas de trabajo. \cite{huang2026} abordaron y resolvieron este problema directamente mediante la estratégica fusión de señales de vibración y corriente. Esta técnica se complementó excepcionalmente con un esquema avanzado de algoritmos de aprendizaje por transferencia completamente supervisado.

Su estrategia combina de manera sumamente eficiente el preentrenamiento algorítmico y un riguroso ajuste inicial de parámetros. Aunque la metodología requiere obtener algunos datos etiquetados nuevos, minimiza de forma muy importante la valiosa información necesaria para lograr un reentrenamiento efectivo. La calidad del procesamiento final de las señales constituye otro aspecto fundamental, crítico y determinante en la fiabilidad.

Las señales medidas operativamente en una máquina CNC pueden contener elevado ruido eléctrico, indeseables vibraciones estructurales y cambios de carga sumamente severos. Por esta precisa razón, resulta estrictamente necesario implementar siempre procedimientos eficaces de filtrado digital y segmentación temporal avanzada. En este ámbito técnico, \cite{jia2025} propusieron un método de reducción de ruido basado eficientemente en descomposición modal variacional adaptativa.

De una manera bastante similar, el detallado estudio sobre la técnica ICFO-SVMD propuso una prometedora metodología que integra optimización algorítmica y descomposición modal \citep{icfo2026}. El principal objetivo analítico fue lograr procesar señales marcadamente no estacionarias que estaban severamente afectadas por ruido electromagnético de origen industrial. En consecuencia, el presente proyecto aplicará un esquema de procesamiento robusto y fuertemente escalonado por etapas.

Además de requerir una constante mejora algorítmica, la incorporación de novedosa conectividad IoT logra potenciar sustancialmente las actuales capacidades del monitoreo. Permite fundamentalmente que el sistema embebido se convierta rápidamente en una moderna plataforma remota de apoyo a la toma de decisiones. El interesante caso de estudio aplicado de la prestigiosa Universidad de Sheffield se erige como un antecedente sumamente relevante \citep{sheffield2021}.

Finalmente, \cite{hamasha2025} propusieron un complejo y completo marco integral tecnológico e industrial de indudable y gran envergadura. Dicho marco funcional articula eficientemente diversos sensores IoT, procesamiento local en el extremo o borde operativo (\textit{Edge AI}), analítica predictiva y monitoreo remoto. El estudio logra identificar cuantiosos beneficios clave como la reducción drástica de latencia y tiempos de respuesta, advirtiendo enormes desafíos en ciberseguridad industrial.

A partir de todos los numerosos antecedentes exhaustivamente revisados, se observa que ya existen sólidas investigaciones consolidadas relativas a detección y fallas. Sin embargo, dichos estudios literarios abordan los diferentes componentes tecnológicos mayoritariamente de manera totalmente independiente y sumamente aislada. Por consiguiente, el gran y ambicioso aporte ingenieril de este proyecto radica en la inusual integración simultánea de estas tecnologías operando en el entorno real de MAINSA.
